\subsection{Dualizing modules}

DEFINITION 1.--- Let A be a Noetherian ring. A Λ-module Ω is said to be dualizing if it is finitely generated and if, for every maximal ideal m of Λ, the A/m-vector space Ext \(_{A}^{i}\) (A/m, Ω) is zero for i ≠ ht(m) and of dimension 1 for i = ht(m).

For every maximal ideal m of A and every integer i, the A/m-vector space \(\operatorname{Ext}_{\mathrm{A}}^{i}(\mathrm{A}/\mathfrak{m},\Omega)\) is canonically isomorphic to \(\operatorname{Ext}_{\mathrm{A}_{\mathfrak{m}}}^{i}(\mathrm{A}/\mathfrak{m},\Omega_{\mathfrak{m}})\) By consequence, for a finitely generated A-module \(\Omega\) be dualizing, it is necessary and sufficient that the \(A_{m}\) -module \(\Omega_{m}\) so that it is dualizing for every maximal ideal m of A.

Examples.--- 1) If the ring A is local and Artinian, the dualizing A-modules are the finitely generated injective A-modules. \(\Omega\) such that \(\mathrm{Hom}_{\mathrm{A}}(\kappa_{\mathrm{A}},\Omega)\) Let \( E \) be of dimension 1 ( \(\S\) 3, no. 3, prop. 6), that is, the Matlis A-modules ( \(\S\) 8, no. 3).

\begin{enumerate}
\def\labelenumi{\arabic{enumi})}
\setcounter{enumi}{1}
\tightlist
\item
  For a Noetherian ring A to be Gorenstein, it is necessary and sufficient that the A-module A be dualizing (§ 3, no. 7, prop. 11). In particular, the A-module A is dualizing when A is regular.
\end{enumerate}

Remarks. -- 1) Let A be a Noetherian local ring and \(\Omega\) a finitely generated A-module. The residue field \(\kappa_{\widehat{A}}\) is identified with \(\kappa_{A}\) , and the \(\widehat{A}\) -module \(\widehat{\Omega}\) with \(\widehat{A} \otimes_{A} \Omega\) (III, § 3, no. 4, th. 3). It then follows from A, X, p.~111, prop. 10 that the \(\kappa_{A}\) -vector space \(\mathrm{Ext}_{\mathbf{A}}^{i}(\kappa_{\mathbf{A}}, \Omega)\) is canonically isomorphic to \(\mathrm{Ext}_{\widehat{A}}^{i}(\kappa_{\widehat{A}}, \widehat{\Omega})\) . Consequently, for the \(A\) -module \(\Omega\) to be dualizing, it is necessary and sufficient that the \(\widehat{A}\) -module \(\widehat{\Omega}\) be dualizing.

\begin{enumerate}
\def\labelenumi{\arabic{enumi})}
\setcounter{enumi}{1}
\tightlist
\item
  Let \(\Omega\) a dualizing A-module; for every rank-1 projective A-module L, the A-module \(\Omega\otimes_{\mathrm{A}}\mathrm{L}\) is dualizing (A, X, p.~108, prop. 7, b)). We shall see below ( \(n^{\circ}\) 4, prop. 6) that every dualizing A-module is isomorphic to a module of this form.
\end{enumerate}

PROPOSITION 1.--- Let A be a Noetherian ring and Ω a dualizing A-module.

\begin{enumerate}
\def\labelenumi{\alph{enumi})}
\tightlist
\item
  Then A is a Macaulay ring, and the A-module Ω is Macaulayan.
\end{enumerate}

\[
\text { b) } O n a \mathrm{di} _ {\mathrm{A}} (\Omega) = \dim (\Omega) = \dim (\mathrm{A}).
\]

Suppose first that the ring A is local, and denote \(d\) its dimension. Prop. 6 of § 3, no. 3 implies \(\mathrm{di}_{\Lambda}(\Omega) = d\) , hence \(\mathrm{prof}(\mathrm{A}) = d\) by Prop. 9 of § 3, no. 6, so that A is a Macaulay ring. Moreover, we have \(\mathrm{prof}(\Omega) = d\) by definition of depth; since we have \(\mathrm{prof}(\Omega) \leqslant \dim(\Omega) \leqslant d\) from which we deduce the proposition in this case.

In the general case, the \(A_{m}\) -module \(\Omega_{m}\) is dualizing for every maximal ideal m of A, hence \(A_{m}\) is a Macaulay ring and \(\Omega_{m}\) a \(\Lambda_{m}\) a). Moreover, we have \(\mathrm{di}_{\mathrm{A}_{\mathfrak{m}}(\Omega_{\mathfrak{m}})} = \dim(\Omega_{\mathfrak{m}}) = \dim(\Lambda_{\mathfrak{m}})\) for every maximal ideal m, whence b) by passing to the supremum (§ 3, no. 2, prop. 3).

PROPOSITION 2.--- Let A be a Noetherian ring, Ω a dualizing A-module. For every prime ideal p of A, the A- \(_{p}\) -module Ω \(_{p}\) is dualizing.

Consider a saturated chain \(\mathfrak{p} \subset \mathfrak{p}_1 \subset \ldots \subset \mathfrak{p}_r\) of prime ideals of A such that the ideal \(\mathfrak{p}_r\) so is maximal. Proceeding by induction on \(r\) one may assume that the \(\mathrm{A}_{\mathfrak{p}_1}\) -module \(\Omega_{\mathfrak{p}_1}\) is dualizing. Replacing \(\Lambda\) by \(\mathrm{A}_{\mathfrak{p}_1}\) and \(\mathfrak{p}\) by \(\mathfrak{p}\mathrm{A}_{\mathfrak{p}_1}\) , we reduce to the case where the ring A is local and the chain \(\mathfrak{p} \subset \mathfrak{m}_{\Lambda}\) is saturated.

Let us then set \(d = \dim(A) = \mathrm{ht}(\mathfrak{m}_{\mathbf{A}})\) . One has \(\dim(A_{\mathfrak{p}}) = \mathrm{ht}(\mathfrak{p}) = d - 1\) since A is a Macaulay ring (§ 2, no. 2, cor. of prop. 2). For every integer \(i\) , the \(A_{\mathfrak{p}}\) -module \(\operatorname{Ext}_{A_{\mathfrak{p}}}^{i}(\kappa(\mathfrak{p}), \Omega_{\mathfrak{p}})\) is isomorphic to \(\operatorname{Ext}_{\Lambda}^{i}(A/\mathfrak{p}, \Omega)_{\mathfrak{p}}\) By § 3, no. 2, prop. 2, it therefore suffices to prove that the \(A/\mathfrak{p}\) -module \(\operatorname{Ext}_{A}^{i}(A/\mathfrak{p}, \Omega)\) is zero for \(i \neq d - 1\) and of rank one for \(i = d - 1\) .

Let \(x\) be an element of \(\mathfrak{m}_{\mathrm{A}} - \mathfrak{p}\) , and by \(\overline{x}\) its class in \(A / \mathfrak{p}\) . Consider the exact sequence of A-modules

\[
0 \rightarrow \mathrm{A} / \mathfrak {p} \xrightarrow {\overline {{{x}}}} \mathrm{A} / \mathfrak {p} \longrightarrow \mathrm{A} / (\mathfrak {p} + x \mathrm{A}) \rightarrow 0.
\]

The A-module \(\mathrm{A} / (\mathfrak{p} + x\mathrm{A})\) is of finite length since its support is reduced to \(\mathfrak{m}_{\mathrm{A}}\) ; since the A-module \(\Omega\) is dualizing, we have \(\operatorname{Ext}_{\mathrm{A}}^{i}(\mathrm{A} / (\mathfrak{p} + x\mathrm{A}), \Omega) = 0\) for \(i \neq d\) (§ 8, no. 5, example 3). One then deduces from the exact sequence of extension modules associated with the above sequence and with \(\Omega\) that the homothety of ratio \(x\) in the A-module \(\operatorname{Ext}_{\mathrm{A}}^{i}(\mathrm{A} / \mathfrak{p}, \Omega)\) is surjective for \(i \neq d - 1\) , which implies that this module is zero (Nakayama's lemma). In particular \(\operatorname{Ext}_{\mathrm{A}}^{d}(\mathrm{A} / \mathfrak{p}, \Omega)\) is zero, and one obtains an exact sequence

\[
0 \rightarrow \operatorname{Ext} _ {\Lambda} ^ {d - 1} (\mathrm{A} / \mathfrak {p}, \Omega) \xrightarrow {x} \operatorname{Ext} _ {\Lambda} ^ {d - 1} (\mathrm{A} / \mathfrak {p}, \Omega) \longrightarrow \operatorname{Ext} _ {\Lambda} ^ {d} (\mathrm{A} / (\mathfrak {p} + x \Lambda), \Omega) \rightarrow 0.
\]

We have \(\operatorname{long}_{\mathrm{A}}(\operatorname{Ext}_{\mathrm{A}}^{i}(\mathrm{A}/(\mathfrak{p}+x\mathrm{A}),\Omega))=\operatorname{long}_{\mathrm{A}}(\mathrm{A}/(\mathfrak{p}+x\mathrm{A}))\) (loc. cit.); the proposition then follows from the following lemma, applied to the ring B = A/p and to the B-module \(M=\operatorname{Ext}_{\Lambda}^{d-1}(\mathrm{A}/\mathfrak{p},\Omega)\) :

Lemma 1. Let B be a local, integral, noetherian ring of dimension 1, and let M be a finitely generated torsion-free B-module. Suppose that one has long \(_{B}\) (M/xM) = long \(_{B}\) (B/xB) for every nonzero element x of B. Then the B-module M has rank 1.

Indeed, let \(r\) the rank of M; there exists a submodule L of M free of rank \(r\) such that M/L is a torsion module (VII, § 4, no. 1, cor. of prop. 1), hence of finite length (VII, § 2, no. 5, lemma 1). The annihilator of M/L is not reduced to 0, and therefore contains a nonzero element \(x\) of \(\mathfrak{m}_{\mathbb{R}}\) . Consider the diagram

commutative

\[
\begin{array}{c c c c c c c} 0 \to & \mathrm{L} & \longrightarrow & \mathrm{M} & \longrightarrow & \mathrm{M} / \mathrm{L} & \to 0 \\ & \Biggl \downarrow x _ {\mathrm{L}} & & \Biggl \downarrow x _ {\mathrm{M}} & & \Biggl \downarrow 0 \\ 0 \to & \mathrm{L} & \longrightarrow & \mathrm{M} & \longrightarrow & \mathrm{M} / \mathrm{L} & \to 0. \end{array}
\]

By the snake lemma (A, X, p.~4, prop. 2), one deduces an exact sequence

\[
0 \rightarrow \mathrm{M} / \mathrm{L} \longrightarrow \mathrm{L} / x \mathrm{L} \longrightarrow \mathrm{M} / x \mathrm{M} \longrightarrow \mathrm{M} / \mathrm{L} \rightarrow 0,
\]

whence long(M/xM) = long(L/xL). Since long(M/xM) = long(B/xB) by hypothesis and long(L/xL) = r long(B/xB), one deduces r = 1.

COROLLARY 1. For every multiplicative subset S of A, the \(\mathrm{S}^{-1}\mathrm{A}\) -module \(\mathrm{S}^{-1}\Omega\) is dualizing.

COROLLARY 2.--- The support of Ω is equal to Spec(A).

Indeed a dualizing module over a local ring is nonzero by definition.

COROLLARY 3. Let M be a finitely generated A-module, and let i be an integer. The A-module \(\operatorname{Ext}_{\mathrm{A}}^{i}(\mathrm{M},\Omega)\) is finitely generated, and its support has codimension \(\geqslant i\) in \(\operatorname{Spec}(\mathrm{A})\) .

The first assertion follows from A, X, p.~108, cor. Let \(\mathfrak{p}\) a prime ideal of the support of \(\operatorname{Ext}_{\mathrm{A}}^{i}(\mathrm{M},\Omega)\) . One has \(\operatorname{Ext}_{\mathrm{A}}^{i}(\mathrm{M},\Omega)_{\mathfrak{p}} \neq 0\) , hence \(\operatorname{Ext}_{\mathrm{A}_{\mathfrak{p}}}^{i}(\mathrm{M}_{\mathfrak{p}},\Omega_{\mathfrak{p}}) \neq 0\) ( \(\S 3, \mathrm{n}^{\circ}2\) , prop. 2), which implies \(\mathrm{di}_{\mathrm{A}_{\mathfrak{p}}}(\Omega_{\mathfrak{p}}) \geqslant i\) . Since \(\Omega_{\mathfrak{p}}\) is a \(\mathrm{A}_{\mathfrak{p}}\) By the dualizing module (prop. 2), we have \(\mathrm{di}_{\mathrm{A}_{\mathfrak{p}}}(\Omega_{\mathfrak{p}}) = \dim(\mathrm{A}_{\mathfrak{p}})\) (prop. 1), whence the corollary.

PROPOSITION 3.--- Let A be a Noetherian local ring, Ω a dualizing A-module, and M a finitely generated A-module.

\[
\text { a) } O n a \operatorname{Ext} _ {\mathrm{A}} ^ {i} (\mathrm{M}, \Omega) = 0 \text { for } i <   \dim (\mathrm{A}) \dots \dim_ {\mathrm{A}} (\mathrm{M}).
\]

\begin{enumerate}
\def\labelenumi{\alph{enumi})}
\setcounter{enumi}{1}
\tightlist
\item
  Set \(c = \dim(\mathrm{A}) - \dim_{\mathrm{A}}(\mathrm{M})\) If M is nonzero, the A-module \(\operatorname{Ext}_{\mathrm{A}}^{c}(\mathrm{M}, \Omega)\) is not zero.
\end{enumerate}

\[
\mathrm{c)} O n a \operatorname{Ext} _ {\mathrm{A}} ^ {i} (\mathrm{M}, \Omega) = 0 p o u r i > \dim (\mathrm{A}) - \operatorname{prof} _ {\mathrm{A}} (\mathrm{M}).
\]

Suppose M is nonzero and denote by F its support. By prop. 9 of § 1, no. 5, the conjunction of assertions a) and b) is equivalent to \(\operatorname{prof}_{\mathbb{F}}(\Omega)=c\) Now since \(\Omega\) is Macaulaycn and that its support is equal to \(\operatorname{Spec}(\mathbb{A})\) (prop. 1 and cor. 2 of prop. 2), we have

\[
\operatorname{prof} _ {\mathrm{F}} (\Omega) = \operatorname{codim} (\mathrm{F}, \operatorname{Spec} (\Lambda)) = c
\]

(§ 2, no. 1, cor. of prop. 1 and no. 2, cor. of prop. 2).

Let us prove c) by induction on the depth of M. If \(\mathrm{prof}_{\Lambda}(\mathbf{M}) = 0\) , we indeed have \(\mathrm{Ext}_{\Lambda}^{i}(\mathbf{M},\Omega) = 0\) for \(i > \dim (\mathbf{A})\) , since \(\mathrm{di}_{\Lambda}(\Omega) = \dim (\mathbf{A})\) (prop. 1). Suppose \(\mathrm{prof}_{\Lambda}(\mathbf{M}) > 0\) then there exists an element \(x\) of \(\mathfrak{m}_{\mathbf{A}}\) such that the homothety of ratio \(x\) injective in M. We have \(\mathrm{prof}_{\Lambda}(\mathbf{M} / x\mathbf{M}) = \mathrm{prof}_{\Lambda}(\mathbf{M}) - 1\) ( \(\S 1, \mathrm{n}^{\circ}4\) prop. 7).

Consider the exact sequence of extension modules

\[
\operatorname{Ext} _ {\mathrm{A}} ^ {i} (\mathrm{M}, \Omega) \xrightarrow {x} \operatorname{Ext} _ {\mathrm{A}} ^ {i} (\mathrm{M}, \Omega) \longrightarrow \operatorname{Ext} _ {\Lambda} ^ {i + 1} (\mathrm{M} / x \mathrm{M}, \Omega)
\]

associated with the exact sequence

\[
0 \to \mathrm{M} \xrightarrow {x} \mathrm{M} \longrightarrow \mathrm{M} / x \mathrm{M} \to 0.
\]

For \(i > \dim(\mathrm{A}) - \operatorname{prof}_{\mathrm{A}}(\mathrm{M})\) , the A-module \(\operatorname{Ext}_{\mathrm{A}}^{i+1}(\mathrm{M}/x\mathrm{M},\Omega)\) is zero by the induction hypothesis, hence the homothety of ratio \(x\) is surjective in \(\operatorname{Ext}_{\mathrm{A}}^{i}(\mathrm{M},\Omega)\) , which implies that this A-module is zero (Nakayama's lemma). This proves c).

COROLLARY.--- If M is Macaulay, one has Ext \(^{i}\) A(M, Ω) = 0 for i ≠ c; the A-module Ext \(^{c}\) A(M, Ω) is Macaulay, and its support is equal to that of M.

The first assertion follows from prop. 3, a) and c). Let \(\mathfrak{p} \in \operatorname{Supp}(M)\) ; by prop. 1 of § 2, no. 1, applied to M and to A, we have

\[
\dim (A _ {p}) - \dim_ {A _ {p}} (M _ {p}) = \dim (A) - \dim_ {A} (M) = c;
\]

since the \(A_{p}\) -module \(\Omega_{p}\) is dualizing (prop. 2), it follows from prop. 3, b) that the \(A_{p}\) -module \(\operatorname{Ext}_{A_{p}}^{c}(M_{p},\Omega_{p})\) is not zero. Consequently, the support of \(\operatorname{Ext}_{A}^{c}(M,\Omega)\) is equal to that of M.

Let us finally prove, by induction on \(\dim(M)\) that the A-module \(\operatorname{Ext}_{A}^{c}(M,\Omega)\) is Macaulay. The assertion is satisfied when \(\dim(M)=0\) since every module of finite length is Macaulay. Suppose \(\dim(M)>0\) and let us choose an element \(x\) of \(\mathfrak{m}_{\mathrm{A}}\) such that the homothety \(x_{\mathrm{M}}\) so that it is injective. The A-module M/ \(x\mathrm{M}\) is Macaulay \(\S 2\) (no. 3, prop. 4), of dimension \(\dim(M)-1\) . In view of the above, the exact sequence of extension modules associated with the exact sequence \(0\to\mathrm{M}\xrightarrow{x}\mathrm{M}\longrightarrow\mathrm{M}/x\mathrm{M}\to0\) reduces to

\[
0 \to \operatorname{Ext} _ {\mathrm{A}} ^ {c} (\mathrm{M}, \Omega) \xrightarrow {x} \operatorname{Ext} _ {\mathrm{A}} ^ {c} (\mathrm{M}, \Omega) \longrightarrow \operatorname{Ext} _ {\mathrm{A}} ^ {c + 1} (\mathrm{M} / x \mathrm{M}, \Omega) \to 0;
\]

Proposition 4 of §2, No. 3 and the induction hypothesis then imply that Ext \(_{A}^{c}\) (M, Ω) is Macaulay, whence the corollary.

